\section{The cubic exponent: an explicit Eisenstein descent}
\label{sec:three}

We prove the stronger signed statement
\[
  xyz\ne0\quad\Longrightarrow\quad x^3+y^3+z^3\ne0
  \qquad(x,y,z\in\mathbb Z).
\]
Signed integers here are finite codes for a sign and an absolute value.
In particular, neither complex numbers nor an infinite number field is
needed.  The descent decreases the absolute value of the coordinate
divisible by~$3$.  Its construction and its termination both have
primitive-recursive bounds.

\subsection{The finite arithmetic needed below}

For clarity, we record the parts of Lemma~\ref{lem:integer-tools} used
below, in their algorithmic forms.  Trial division through $2,\ldots,|m|$
produces a prime factorization of each nonzero $m$.  Euclid's algorithm
produces $\gcd(u,v)$ and Bezout coefficients.  The usual proofs use division
with remainder and the identity preserving the gcd at one Euclidean
step.  They consequently use induction only over explicitly finite
computations.  Comparing prime exponents in the resulting lists proves:

\begin{enumerate}
\item If $d^3\mid m^3$, with $d>0$, then $d\mid m$.
\item If $u,v,w$ are pairwise coprime signed integers and $uvw$ is a cube,
then each of $u,v,w$ is a signed cube.  Indeed, a prime occurs in at most
one of their absolute values, so its exponent there is a multiple of
three; attach the original sign after taking the positive cube root.
\item If $\gcd(u,v)=1$ and $u^3+v^3+w^3=0$, the three coordinates are
pairwise coprime.  A prime dividing any two divides the third.
\end{enumerate}

These facts do not require an unbounded search for a factor or a root.
Divisors of $m$ are searched below $|m|$, and a cube root of a nonzero
integer $m$ is searched in $[-|m|,|m|]$.  Later applications always supply
the relevant nonzero integers.

\subsection{The pair ring and its Euclidean algorithm}

Let $E$ be the set of pairs of signed integers, written $a+b\omega$, with
\[
 \omega^2+\omega+1=0.
\]
This is a notation for the following explicit operations:
\begin{align*}
 (a,b)+(c,d)&=(a+c,b+d),\\
 (a,b)(c,d)&=(ac-bd,\ ad+bc-bd),\\
 \overline{(a,b)}&=(a-b,-b),\\
 N(a,b)&=a^2-ab+b^2.
\end{align*}
Direct multiplication verifies
\[
 \alpha\bar\alpha=N(\alpha),\qquad
 N(\alpha\beta)=N(\alpha)N(\beta),\qquad
 2N(a,b)=a^2+b^2+(a-b)^2.
\]
Thus $N\ge0$, and $N(\alpha)=0$ exactly when $\alpha=0$.
In particular, this pair ring is an integral domain: if
$\alpha\beta=0$, multiplication of natural-number norms gives
$N(\alpha)N(\beta)=0$, and decidable arithmetic gives $\alpha=0$ or
$\beta=0$.  The last displayed identity also gives the useful bound
\begin{equation}
 |a|,\ |b|,\ |a-b|\ \le 2N(a,b).
 \label{eq:cubic-coordinate-bound}
\end{equation}
Here we used $|t|\le t^2$ for every signed integer $t$.

We next give the division algorithm.  If $\beta=c+d\omega\ne0$, set
$D=N(\beta)>0$.  In the rational pair ring the quotient $\alpha/\beta$,
for $\alpha=a+b\omega$, has coordinates
\[
 u=\frac{ac-ad+bd}{D},\qquad v=\frac{bc-ad}{D}.
\]
Choose integers $q_0,q_1$ with $|u-q_0|,|v-q_1|\le1/2$, using signed
integer division and a fixed tie convention.  Put
$q=q_0+q_1\omega$ and $r=\alpha-q\beta$.  For rational numbers
$s,t$ of absolute value at most $1/2$,
\[
 s^2-st+t^2\le\tfrac14+\tfrac14+\tfrac14=\tfrac34.
\]
Consequently
\[
 N(r)\le\tfrac34N(\beta)<N(\beta).
\]
The proof is entirely rational arithmetic; denominators can simply be
cleared.  Thus Euclid's algorithm terminates after at most
$N(\beta)+1$ divisions, including a possible initial exchange.  Repeated
substitution in the division equations gives a gcd $g$ and a Bezout
identity $g=s\alpha+t\beta$.

Here are the precise consequences of this algorithm, with their proofs.
An element is a unit exactly when its norm is~$1$: necessity follows
from multiplicativity, and sufficiency from $\alpha\bar\alpha=1$.
The six units are
\[
 \{1,-1,\omega,-\omega,\omega^2,-\omega^2\}.
\]
To check completeness of this list, write
$4N(a,b)=(2a-b)^2+3b^2$.  When $N=1$, one has
$b\in\{-1,0,1\}$, and checking these three possibilities gives exactly
the six displayed elements.

Every nonzero nonunit has a factorization into irreducibles.  To construct
one, search for a nonunit proper divisor in the box
$|a|,|b|\le2N(\alpha)$, and split when one is found.  A divisor of
$\alpha$ has norm at most $N(\alpha)$, hence lies in this box by
\eqref{eq:cubic-coordinate-bound}; divisibility is checked by the explicit
quotient formula above.  If no such divisor exists, the element is
irreducible.  A proper splitting replaces a norm by two smaller positive
norms, both at least two.  Equivalently one may repeatedly split a
factor list; after $k$ proper splittings its product has at least $k+1$
nonunit factors, so $2^{k+1}\le N(\alpha)$.  In particular a loop with
$N(\alpha)$ rounds suffices.

Every irreducible is prime.  If an irreducible $\pi$ does not divide
$\eta$, its gcd with $\eta$ is a unit.  The computed Bezout identity
can then be normalized to $u\pi+v\eta=1$.  If $\pi\mid\eta\theta$,
multiplying this identity by $\theta$ proves $\pi\mid\theta$.
Finally, factorization is unique up to order and units: compare an
irreducible in one factorization with the product in the other, use
primality to find an associate, cancel it, and repeat through the finite
factor lists.  This is the only unique-factorization result used below;
it has just been proved for our explicit pair ring.

\subsection{The prime above three}

Put $\lambda=1-\omega=(1,-1)$.  Calculation gives
\[
 N(\lambda)=3,\qquad \lambda^2=-3\omega,
 \qquad \bar\lambda=-\omega^2\lambda.
\]
The norm $3$ is an ordinary prime, so any nonunit proper divisor of
$\lambda$ would have an impossible intermediate norm.  Thus $\lambda$
is irreducible and therefore prime.  For $\alpha=a+b\omega$,
\begin{equation}
 \lambda\mid\alpha\ \Longleftrightarrow\ 3\mid a+b,
 \qquad
 \lambda^2\mid\alpha\ \Longleftrightarrow\ (3\mid a\ \text{and}\ 3\mid b).
 \label{eq:cubic-lambda}
\end{equation}
The first equivalence follows from the quotient formula
\[
 \frac{a+b\omega}{1-\omega}
 =\frac{2a-b}{3}+\frac{a+b}{3}\omega:
 \quad 2a-b=2(a+b)-3b.
\]
The second follows from $\lambda^2=-3\omega$ and the fact that
$\omega$ is a unit.  These observations also prove, without any
abstract quotient-ring construction, all ramification facts needed here.

\subsection{Normalization of a hypothetical solution}

Suppose $x,y,z$ are nonzero signed integers with
$x^3+y^3+z^3=0$.  Divide by $g=\gcd(|x|,|y|)>0$.
Since $g^3\mid z^3$, the elementary cube-divisibility fact gives
$g\mid z$.  The resulting integer triple is a solution with
$\gcd(x,y)=1$, and hence is pairwise coprime.  We keep the same letters
for this primitive triple.

The cubes modulo~$9$ are $0,1,-1$.  If none of $x,y,z$ is divisible
by~$3$, the sum of their cubes is one of $-3,-1,1,3$ modulo~$9$, never
zero.  Thus one coordinate is divisible by~$3$, and pairwise coprimality
shows that exactly one is.  Permute the coordinates to arrange
\begin{equation}
 3\mid z,\qquad 3\nmid xy,\qquad 3\mid x+y.
 \label{eq:cubic-normalization}
\end{equation}
In fact $9\mid x+y$.  Write $x+y=3k$.  Then
\[
 -z^3=x^3+y^3=(x+y)((x+y)^2-3xy)
                  =9k(3k^2-xy).
\]
The left side is divisible by~$27$, while $3k^2-xy$ is not divisible
by~$3$.  Cancellation of $9$ and primality of~$3$ give $3\mid k$.
Also $x+y\ne0$, since otherwise the original equation would give $z=0$.

\subsection{Extracting the cube in the pair ring}

Set $\alpha=x+y\omega$.  It is nonzero, and
\begin{equation}
 (x+y)(x+y\omega)(x+y\omega^2)=-z^3.
 \label{eq:cubic-factorization}
\end{equation}
By \eqref{eq:cubic-lambda} and
\eqref{eq:cubic-normalization}, $\lambda$ divides $\alpha$ exactly once:
$\lambda\mid\alpha$ but $\lambda^2\nmid\alpha$.

Any irreducible $\pi$ not associate to $\lambda$ that divides $\alpha$
divides neither of the other two factors on the left of
\eqref{eq:cubic-factorization}.  Indeed, a common divisor of $\alpha$
and $x+y$ divides their difference $-y\lambda$.  A common divisor of
$\alpha$ and $\bar\alpha=x+y\omega^2$ divides
$y(\omega-\omega^2)=y\omega\lambda$.  In either case primality and
$\pi\nmid\lambda$ give $\pi\mid y$; subtracting $y\omega$ from
$\alpha$ gives $\pi\mid x$.  Ordinary Bezout coefficients for
$\gcd(x,y)=1$ would then make $\pi$ divide~$1$, a contradiction.

Compare the finite irreducible factorizations of both sides of
\eqref{eq:cubic-factorization}.  For every irreducible other than
$\lambda$, its exponent in $\alpha$ is its entire exponent on the
left, and so is a multiple of three.  The exponent of $\lambda$ in
$\alpha$ is one.  Therefore there are a unit $\varepsilon$ and an
element $\gamma=a+b\omega$ such that
\begin{equation}
 x+y\omega=\varepsilon\lambda(a+b\omega)^3,
 \qquad \lambda\nmid(a+b\omega).
 \label{eq:cubic-extraction}
\end{equation}
This is a constructive extraction from a factor list.  Alternatively,
one can find the same witness by a bounded search: taking norms shows
$3N(\gamma)^3=N(\alpha)$, so
$|a|,|b|\le2N(\alpha)$; search this box and the six units.

\subsection{The unit calculation, including the exceptional possibilities}

Write
\[
 A=a^3-3ab^2+b^3,\qquad B=3ab(a-b).
\]
Direct multiplication gives
\[
 (a+b\omega)^3=A+B\omega,
 \qquad
 \lambda(a+b\omega)^3=(A+B)+(2B-A)\omega.
\]
The sums of the two coefficients after multiplication by the three
units $1,\omega,\omega^2$ are, respectively,
\[
 3B,\qquad 3(A-B),\qquad -3A.
\]
Multiplication by their negatives just changes the signs.  Now
$3\mid B$, while
\[
 A\equiv a+b\not\equiv0\pmod3;
\]
the final inequality is exactly $\lambda\nmid\gamma$ from
\eqref{eq:cubic-lambda}.  Since $9\mid x+y$, the choices
$\varepsilon=\pm\omega,\pm\omega^2$ are impossible.  Consequently
\begin{equation}
 \varepsilon\in\{1,-1\},\qquad
 x+y=\varepsilon\,9ab(a-b).
 \label{eq:cubic-unit}
\end{equation}
This check is essential: discarding units without this argument would
leave a gap in the descent.

\subsection{The smaller solution}

Put $M=N(a,b)>0$.  Norms in \eqref{eq:cubic-extraction} give
$N(x+y\omega)=3M^3$.  Together with
\eqref{eq:cubic-unit}, this gives
\begin{equation}
 -z^3=(x+y)N(x+y\omega)
      =\varepsilon\,27M^3ab(a-b).
 \label{eq:cubic-product}
\end{equation}
Thus $(3M)^3\mid z^3$, and hence $3M\mid z$.  Set
$h=z/(3M)$.  Equation \eqref{eq:cubic-product} now says
\[
 ab(a-b)=-\varepsilon h^3=(-\varepsilon h)^3.
\]
None of $a,b,a-b,h$ is zero, since $x+y$ and $z$ are nonzero.

The integers $a,b$ are coprime.  If a positive integer $d$ divides both,
then $d^3$ divides both coordinates of
$\varepsilon\lambda(a+b\omega)^3=x+y\omega$; hence $d^3$ divides
$\gcd(x,y)=1$, so $d=1$.  It follows that $a,b,a-b$ are pairwise
coprime.  The ordinary integer cube-splitting lemma therefore supplies
nonzero signed integers $r,s,t$ with
\[
 a=r^3,\qquad b=s^3,\qquad a-b=t^3.
\]
We have obtained another cubic solution,
\[
 s^3+t^3+(-r)^3=0.
\]
It is primitive: any prime dividing two of $r,s,t$ would divide two of
$a,b,a-b$.  By the same modulo-$9$ calculation as before, exactly one
of $s,t,-r$ is divisible by~$3$.  Permute these three entries so that
this coordinate is called $z'$.

The decrease is explicit.  Since $|u|\le|u|^3$ for nonzero integers,
\[
 |r|\le|a|,\qquad |s|\le|b|,\qquad |t|\le|a-b|.
\]
By \eqref{eq:cubic-coordinate-bound} each right side is at most $2M$.
On the other hand, $z=3Mh$ with $|h|\ge1$, so
\begin{equation}
 0<|z'|\le2M<3M\le|z|.
 \label{eq:cubic-descent}
\end{equation}
We have therefore constructed, from every normalized primitive
solution, another such solution with a smaller positive distinguished
coordinate.

\subsection{Finite termination in iPRA}

It remains to instantiate Lemma~\ref{lem:clocked} explicitly, rather
than appeal to an infinite descending sequence or to a least
counterexample.  Let $V(S)$ be the
primitive-recursive Boolean predicate saying that a coded triple $S$
has nonzero pairwise coprime signed coordinates, their cubes sum to
zero, and the third coordinate alone is divisible by~$3$.  Let
$m(S)$ be the absolute value of this third coordinate.  The algorithms
above define a total primitive-recursive function $D$: when $V(S)$ is
true, perform the bounded cube extraction and subsequent bounded
integer root searches, and return the first resulting normalized
triple in a fixed order; otherwise return a fixed dummy triple.  The
proof above establishes
\[
 V(S)=1\quad\Longrightarrow\quad
 V(D(S))=1\ \text{ and }\ 0<m(D(S))<m(S).
\]
Define $S_0=S$ and $S_{k+1}=D(S_k)$ by primitive recursion.  If
$m=m(S)$ and $V(S)=1$, quantifier-free induction on $k$, with the
guard $k\le m$, proves
\[
 k\le m\quad\Longrightarrow\quad V(S_k)=1\ \text{ and }\ m(S_k)+k\le m.
\]
At $k=m$ this gives $m(S_m)\le0$, whereas $V(S_m)=1$ gives
$m(S_m)>0$.  This is a contradiction in decidable natural-number
arithmetic.  The induction formula contains only primitive-recursive
terms and decidable relations; bounded conjunctions, if used to check
factor lists, are represented by Boolean folds.  No excluded middle
for arbitrary formulas and no unbounded minimization have been used.

For an explicit size audit, if both Euclidean inputs have norm at most
$H\ge1$, their coordinates are bounded by $2H$.  In the quotient
formula, the numerators are bounded in absolute value by $12H^2$,
so rounded quotient coordinates are bounded by $Q=12H^2+1$.
There are at most $H+2$ steps.  Bezout coefficients are therefore
bounded, for example, by
\[
 (3Q+2)^{H+3}.
\]
Indeed a coefficient update is an old coefficient minus a product of
a quotient and another coefficient, and pair multiplication increases
the coordinate bound by at most a factor of $3Q$.
Factorization uses at most $N(\alpha)$ factors and rounds; each
candidate factor has coordinates bounded by $2N(\alpha)$.
The cube-extraction box has the same bound.  These coarse bounds are
all primitive recursive and suffice for coding every invariant by
bounded searches and finite lists.  The descent itself takes only
$m(S)$ steps; by \eqref{eq:cubic-descent}, after its first step all
three coordinate absolute values are below the original $m(S)$.

We have proved the signed cubic statement, and therefore
\[
 \forall x,y,z\in\mathbb N_{>0},\qquad x^3+y^3\ne z^3.
\]
Replacing $x,y,z$ by $x^k,y^k,z^k$ gives FLT for every positive
multiple $3k$ of~$3$.  This final substitution is uniform in~$k$,
since exponentiation itself is primitive recursive.

\subsection{Lean obligations for this proof}

A direct formalization should use a pair type with signed-integer
coordinates, rather than import a general number-field construction.
The following are the exact successive obligations.

\begin{enumerate}
\item Define pair addition, multiplication, conjugation and norm by the
four displayed formulas.  Prove their ring identities and norm
multiplicativity by polynomial normalization, and prove
\eqref{eq:cubic-coordinate-bound} by integer inequalities.
\item Implement signed nearest-integer division.  Prove the Euclidean
remainder bound after clearing denominators.  Define the gcd computation
with explicit norm fuel and prove its Bezout invariant.
\item Implement bounded irreducible factorization.  Prove primality by
the Bezout identity and uniqueness by finite cancellation.  Formalize
the six-unit table and \eqref{eq:cubic-lambda} as coordinate lemmas.
\item Define the Boolean predicate $V$, the bounded witness extractor
in \eqref{eq:cubic-extraction}, and the descent function $D$.  The
core theorem is the concrete implication
$V(S)=1\Rightarrow V(D(S))=1\land m(D(S))<m(S)$.
The unit table and polynomial identities above give its arithmetic
part; integer coprime cube extraction gives its final three roots.
\item Prove the guarded finite-iteration invariant and derive FLT for
$3$, then for $3k$ by the exponent law $(x^k)^3=x^{3k}$.
\item Separately translate these bounded programs, their defining
equations and their quantifier-free induction instances into the
chosen object-language presentation of iPRA.  A Lean proof about
ordinary integers alone does not perform this last step automatically.
\end{enumerate}

The existing Mathlib declaration
\texttt{fermatLastTheoremThree} in
\nolinkurl{Mathlib.NumberTheory.FLT.Three} is a useful independent
semantic comparison.  Its cyclotomic-field imports are not used as an
unexamined certificate of iPRA-provability in the argument above.
